\chapter{Esercitazione E4: collegamento, fisico, sicurezza e anonimato}
\titolobreve{Esercitazione E4}

L'ultimo pacchetto: livello di collegamento e fisico, sicurezza e anonimato, e alla fine l'intera pila insieme, osservata sulla chat che
avete scritto. L'esercizio $n$ di questo capitolo è l'\textbf{E4.$n$} delle slide. Quasi tutti sono su carta; tre chiedono di catturare
sulla propria macchina, gli ultimi due di eseguire la propria chat. Dove un esercizio dà un codice, una chiave o una tabella, dà tutto ciò
che serve.

\info{Strumenti}{
    Wireshark sull'interfaccia che si usa davvero (il loopback non mostra né Ethernet né ARP né DHCP), e da terminale \texttt{ip},
    \texttt{ping}, \texttt{curl}, \texttt{nmcli}, \texttt{resolvectl}: \texttt{sudo apt install iproute2 iputils-ping curl}. I comandi con
    \texttt{sudo} cambiano lo stato della macchina: si eseguono dove si è autorizzati. Gli E4.29 ed E4.30 richiedono due macchine sulla
    stessa LAN (una macchina virtuale in modalità bridge conta come seconda macchina).
}

% ══════════════════════════════════════════════════════════════
\section{Framing ed errori}
% ══════════════════════════════════════════════════════════════

\begin{esercizio}{Byte stuffing}
I frame sono delimitati da un byte FLAG, ed ESC protegge un FLAG o un ESC nel payload, come nel corso.
\begin{enumerate}[label=(\alph*)]
    \item Applicare lo stuffing al payload \texttt{A FLAG B ESC C ESC FLAG D}.
    \item Tra due flag il ricevente trova \texttt{A ESC ESC B ESC FLAG C ESC ESC ESC FLAG D}. Che payload consegna?
    \item Un errore su un bit trasforma uno degli ESC di (b) in un FLAG. Che cosa perde il ricevente, e da dove ricomincia?
\end{enumerate}

\svolgimento
\begin{enumerate}[label=(\alph*)]
    \item \texttt{FLAG A ESC FLAG B ESC ESC C ESC ESC ESC FLAG D FLAG}.
    \item \texttt{A ESC B FLAG C ESC FLAG D}: si toglie il primo ESC di ogni coppia e si prende alla lettera il byte che segue.
    \item Il FLAG spurio chiude il frame a metà: il pezzo prima e il pezzo dopo (letto come un frame nuovo, fino al FLAG vero) falliscono
          entrambi il CRC e vengono scartati. Si perde questo frame, e il ricevente si risincronizza al FLAG che lo chiudeva: il frame
          successivo arriva intatto.
\end{enumerate}
\end{esercizio}

\begin{esercizio}{Bit stuffing}
Il flag è \texttt{01111110}.
\begin{enumerate}[label=(\alph*)]
    \item Applicare lo stuffing ai dati \texttt{0111111011111100111110}.
    \item Il ricevente riceve \texttt{0111111010111110101111101111001111110}. Trovare il frame e fare il destuffing del payload.
    \item Perché basta uno 0 dopo cinque 1 a mantenere il flag unico?
    \item Che cosa deve concludere il ricevente se vede sette 1 di seguito?
\end{enumerate}

\svolgimento
\begin{enumerate}[label=(\alph*)]
    \item \texttt{011111\textcolor{netred}{0}1011111\textcolor{netred}{0}10011111\textcolor{netred}{0}0}.
    \item Il flag sta ai bit 0--7 e 29--36; il payload è \texttt{101111101011111011110}. Togliendo gli 0 che seguono cinque 1 (due) resta
          \texttt{1011111101111111110}.
    \item Dopo lo stuffing nei dati non possono mai comparire sei 1 di seguito, e il flag ne contiene sei: la sequenza è riservata ai
          delimitatori.
    \item Che c'è un errore: sette 1 non possono essere né dati né flag. Il frame si scarta (in HDLC è la sequenza di \emph{abort}).
\end{enumerate}
\end{esercizio}

\begin{esercizio}{Parità bidimensionale}
L'informazione in un pacchetto è \texttt{1110 0110 1001 0101}, con parità pari. Dare i bit di parità di uno schema bidimensionale, con un
campo di controllo il più corto possibile.

\svolgimento
Con 16 bit la disposizione $4 \times 4$ minimizza $i + j + 1$ (9 bit, contro gli 11 di una $2 \times 8$). Righe \texttt{1110},
\texttt{0110}, \texttt{1001}, \texttt{0101}: parità di riga \texttt{1 0 0 0}; parità di colonna \texttt{0 1 0 0}; bit d'angolo 1.
\end{esercizio}

\begin{esercizio}{CRC}
Il generatore è $G = \texttt{10011}$.
\begin{enumerate}[label=(\alph*)]
    \item Per $D = \texttt{1010101010}$ trovare i bit CRC $R$. Scrivere $G(x)$ e $D(x)$ come polinomi, e verificare che $M(x) = x^r D(x) -
          R(x)$ sia divisibile per $G(x)$.
    \item Trovare $R$ per $D = \texttt{1000100101}$, \texttt{0101101010} e \texttt{0110100011}.
    \item Con $G = \texttt{1001}$, ogni errore su un singolo bit dà un resto non nullo? Ogni numero dispari di errori? Perché?
\end{enumerate}

\svolgimento
\begin{enumerate}[label=(\alph*)]
    \item $G(x) = x^4 + x + 1$, $D(x) = x^9 + x^7 + x^5 + x^3 + x$. Dividendo \texttt{10101010100000} per \texttt{10011} il resto è
          $R = \texttt{0100}$, cioè $R(x) = x^2$. Si trasmette \texttt{1010101010\,0100}, e dividendolo di nuovo per $G$ il resto è 0: è
          $M(x) = x^4 D(x) + R(x)$, divisibile per costruzione.
    \item \texttt{1000}, \texttt{1111}, \texttt{1110}.
    \item $G(x) = x^3 + 1 = (x + 1)(x^2 + x + 1)$. Ha due termini, quindi non divide mai un $E(x) = x^i$: ogni errore singolo dà resto non
          nullo. Ha il fattore $x + 1$, quindi rileva ogni numero dispari di errori ($E(1) = 1$ per un numero dispari di termini, mentre ogni
          multiplo di $x + 1$ vale 0 in 1).
\end{enumerate}
\end{esercizio}

\begin{esercizio}{Distanza di Hamming}
\begin{enumerate}[label=(\alph*)]
    \item Un codice ha quattro parole: \texttt{0000000000}, \texttt{0000011111}, \texttt{1111100000}, \texttt{1111111111}. Qual è la sua
          distanza di Hamming? Quanti errori rileva sempre, e quanti ne corregge? Si riceve \texttt{0100011101}: quale parola è stata
          inviata, assumendo il minor numero di errori?
    \item Rispondere alle stesse domande per il codice \texttt{00000}, \texttt{01011}, \texttt{10101}, \texttt{11110}, e decodificare
          \texttt{11011}, \texttt{11101} e \texttt{01110}.
    \item Un codice deve correggere ogni errore singolo e rilevare ogni errore doppio. Qual è la minima distanza di Hamming che può avere?
\end{enumerate}

\svolgimento
\begin{enumerate}[label=(\alph*)]
    \item Distanza 5: rileva 4 errori, ne corregge 2. \texttt{0100011101} è a distanza 2 da \texttt{0000011111} e almeno 5 dalle altre:
          è stata inviata \texttt{0000011111}.
    \item Le distanze tra le parole sono 3, 3, 4, 4, 3, 3: distanza 3, rileva 2 errori e ne corregge 1. \texttt{11011} $\rightarrow$
          \texttt{01011} (distanza 1), \texttt{11101} $\rightarrow$ \texttt{10101}, \texttt{01110} $\rightarrow$ \texttt{11110}.
    \item Per correggere $t$ errori e rilevarne $s > t$ serve $d \geq t + s + 1$: con $t = 1$ e $s = 2$, distanza \textbf{4} (le parole
          ricevute con due errori restano a metà strada, e non vengono ``corrette'' verso la parola sbagliata).
\end{enumerate}
\end{esercizio}

% ══════════════════════════════════════════════════════════════
\section{Accesso multiplo}
% ══════════════════════════════════════════════════════════════

\begin{esercizio}{Quanto può rendere ALOHA}
$N$ nodi attivi condividono un canale slotted ALOHA; in ogni slot ciascuno trasmette con probabilità $p$. L'efficienza, la frazione di slot
che portano esattamente un frame, è $Np(1-p)^{N-1}$.
\begin{enumerate}[label=(\alph*)]
    \item Trovare la $p$ che la massimizza per un dato $N$.
    \item Con quella $p$, trovare l'efficienza per $N \to \infty$. Suggerimento: $(1 - 1/N)^N \to 1/e$.
    \item In ALOHA puro l'esponente diventa $2(N-1)$. Spiegare perché, e mostrare che l'efficienza massima è allora $1/(2e)$.
    \item Nel corso gli stessi risultati venivano da $S = Ge^{-2G}$ e $S = Ge^{-G}$. A che cosa corrisponde $G$ in questo modello?
\end{enumerate}

\svolgimento
\begin{enumerate}[label=(\alph*)]
    \item La derivata $N(1-p)^{N-2}\bigl[(1-p) - (N-1)p\bigr]$ si annulla per $p^* = 1/N$.
    \item $N \cdot \frac{1}{N}(1 - \frac{1}{N})^{N-1} = (1 - \frac{1}{N})^{N} / (1 - \frac{1}{N}) \to 1/e \approx 0{,}37$.
    \item In ALOHA puro un frame collide con chiunque inizi nel tempo di frame precedente o durante il suo: ogni altro nodo deve tacere per
          \textbf{due} intervalli, da cui $(1-p)^{2(N-1)}$. Massimizzando si trova $p^* = 1/(2N-1)$, e per $N \to \infty$ l'efficienza tende a
          $\frac{1}{2}\, e^{-1} = 1/(2e) \approx 0{,}18$.
    \item $G = Np$: il numero medio di tentativi di trasmissione per slot.
\end{enumerate}
\end{esercizio}

\begin{esercizio}{Binary exponential backoff}
In CSMA/CD, dopo la $n$-esima collisione sullo stesso frame un nodo sceglie $K$ a caso tra $0, 1, \dots, 2^n - 1$ e aspetta $K$ volte 512
tempi di bit.
\begin{enumerate}[label=(\alph*)]
    \item Dopo la quinta collisione, qual è la probabilità che un nodo scelga $K = 4$? Quanto aspetta allora su Ethernet a 10 Mb/s?
    \item Per $K = 100$, quanto aspetta a 100 Mb/s? E a 1 Gb/s, mantenendo i 512 tempi di bit?
\end{enumerate}

\svolgimento
(a) $K \in \{0, \dots, 31\}$: probabilità $1/32$. Aspetta $4 \cdot 512 = 2048$ tempi di bit, cioè $204{,}8\ \mu$s a 10 Mb/s. (b) $51\,200$
tempi di bit: $512\ \mu$s a 100 Mb/s, $51{,}2\ \mu$s a 1 Gb/s.
\end{esercizio}

\begin{esercizio}{La collisione che passa inosservata}
A e B condividono un canale da 10 Mb/s con CSMA/CD; il ritardo di propagazione tra loro è di 325 tempi di bit. A comincia un frame, e B ne
comincia uno prima che A abbia finito. Un frame è almeno di 512 bit più 64 di preambolo, quindi nel caso peggiore A smette di trasmettere a
576 tempi di bit.
\begin{enumerate}[label=(\alph*)]
    \item Nel caso peggiore, quando arriva ad A il segnale di B? A può finire prima di rilevare B?
    \item Esprimere $\tau$ e $2\tau$ in tempi di bit, e dire che cosa significa la risposta per la lunghezza minima del frame.
\end{enumerate}

\svolgimento
(a) Nel caso peggiore B comincia a $t = 324$, un attimo prima che il primo bit di A lo raggiunga; il primo bit di B arriva ad A a $324 + 325
= 649$. Poiché $649 > 576$, A finisce prima di rilevare B, e crede che il frame sia passato senza collisione. (b) $\tau = 325$, $2\tau =
650$ tempi di bit. Il frame minimo (576) è più corto di $2\tau$: questa rete è troppo lunga per quel minimo. Serve un frame di almeno 650
bit, o una rete più corta.
\end{esercizio}

\begin{esercizio}{Nascosti ed esposti}
Cinque stazioni radio stanno su una linea: A a 0 m, B a 80 m, C a 160 m, D a 240 m, E a 320 m. Ognuna raggiunge le stazioni entro 100 m, e
non oltre.
\begin{enumerate}[label=(\alph*)]
    \item A trasmette a B mentre anche C trasmette a B. Che cosa succede in B? Ascoltare il canale prima di trasmettere l'avrebbe evitato?
    \item B sta trasmettendo ad A, e C vuole trasmettere a D. Che cosa fa fare a C il carrier sense? Era necessario?
    \item Ora A trasmette a B con RTS/CTS. Chi sente l'RTS, chi sente il CTS, e chi deve tacere mentre passano i dati? C può trasmettere a D
          nel frattempo? D può trasmettere a E?
\end{enumerate}

\svolgimento
Ognuno sente solo i vicini a 80 m: A--B, B--C, C--D, D--E.
\begin{enumerate}[label=(\alph*)]
    \item I due frame collidono in B. Il carrier sense non l'avrebbe evitato: A e C distano 160 m e non si sentono (terminale nascosto).
    \item C sente B e rimanda (terminale esposto). Non era necessario: la trasmissione di C raggiunge B e D, non A, e A riceverebbe B senza
          disturbo.
    \item L'RTS di A lo sente solo B. Il CTS di B lo sentono A e C: C deve tacere per tutta la durata dei dati, quindi \textbf{non} può
          trasmettere a D (disturberebbe la ricezione in B). D non ha sentito né RTS né CTS e \textbf{può} trasmettere a E: il suo segnale
          raggiunge C ed E, non B.
\end{enumerate}
\end{esercizio}

\begin{esercizio}{Divisione di codice, decodificata}
Quattro stazioni hanno queste sequenze di chip; una stazione invia la propria sequenza per un 1, la sua negazione per uno 0, niente se
tace. Il ricevente fa il prodotto scalare con la sequenza di una stazione e lo divide per 8.
\begin{center}
\small
\begin{tabular}{l*{8}{r}}
\toprule
A & $-1$ & $-1$ & $-1$ & $+1$ & $+1$ & $-1$ & $+1$ & $+1$ \\
B & $-1$ & $-1$ & $+1$ & $-1$ & $+1$ & $+1$ & $+1$ & $-1$ \\
C & $-1$ & $+1$ & $-1$ & $+1$ & $+1$ & $+1$ & $-1$ & $-1$ \\
D & $-1$ & $+1$ & $-1$ & $-1$ & $-1$ & $-1$ & $+1$ & $-1$ \\
\bottomrule
\end{tabular}
\end{center}
\begin{enumerate}[label=(\alph*)]
    \item Verificare che A e B sono ortogonali, e che A con sé stessa dà 1.
    \item Il ricevente riceve $(-1, +1, -3, +1, -1, -3, +1, +1)$. Quali stazioni hanno trasmesso, e quale bit ha inviato ciascuna?
    \item Che cosa arriva se C invia un 1 e D uno 0, e le altre tacciono?
\end{enumerate}

\svolgimento
\begin{enumerate}[label=(\alph*)]
    \item $A \cdot B = (1 + 1 - 1 - 1 + 1 - 1 + 1 - 1)/8 = 0$; $A \cdot A = 8/8 = 1$. (Le quattro sequenze sono a due a due ortogonali.)
    \item Prodotti normalizzati: A $+1$, B $-1$, C $0$, D $+1$. A ha inviato 1, B ha inviato 0, C ha taciuto, D ha inviato 1.
    \item $C - D = (0, 0, 0, +2, +2, +2, -2, 0)$.
\end{enumerate}
\end{esercizio}

% ══════════════════════════════════════════════════════════════
\section{Indirizzi, switch e ARP}
% ══════════════════════════════════════════════════════════════

\begin{esercizio}{Indirizzi, hop per hop}
Tre LAN sono unite da due router. La subnet 1 contiene gli host A e B e un'interfaccia del router R1; la subnet 2 gli host C e D, l'altra
interfaccia di R1 e un'interfaccia del router R2; la subnet 3 gli host E e F e l'altra interfaccia di R2. Si usano \texttt{192.168.1.0/24},
\texttt{192.168.2.0/24} e \texttt{192.168.3.0/24}.
\begin{enumerate}[label=(\alph*)]
    \item Assegnare un indirizzo IP e un indirizzo MAC a ogni interfaccia di A, dei due router e di F.
    \item A invia un datagramma IP a F. Dare gli indirizzi MAC sorgente e destinazione del frame, e gli indirizzi IP sorgente e destinazione
          del datagramma, (i) da A a R1, (ii) da R1 a R2, (iii) da R2 a F.
    \item Quali query ARP servono se ogni tabella ARP parte vuota, e chi invia ciascuna?
\end{enumerate}

\svolgimento
\begin{enumerate}[label=(\alph*)]
    \item Per esempio: A \texttt{192.168.1.10}, \texttt{00-00-00-00-00-0A}; R1 \texttt{192.168.1.1}, \texttt{00-00-00-00-01-01} e
          \texttt{192.168.2.1}, \texttt{00-00-00-00-01-02}; R2 \texttt{192.168.2.2}, \texttt{00-00-00-00-02-01} e \texttt{192.168.3.1},
          \texttt{00-00-00-00-02-02}; F \texttt{192.168.3.20}, \texttt{00-00-00-00-00-0F}.
    \item \begin{tabular}[t]{llll}
          \toprule
           & \textbf{MAC sorg.} & \textbf{MAC dest.} & \textbf{IP} \\
          \midrule
          (i) & A & R1 (\texttt{01-01}) & A $\rightarrow$ F \\
          (ii) & R1 (\texttt{01-02}) & R2 (\texttt{02-01}) & A $\rightarrow$ F \\
          (iii) & R2 (\texttt{02-02}) & F & A $\rightarrow$ F \\
          \bottomrule
          \end{tabular}

          Gli indirizzi IP non cambiano; i MAC cambiano a ogni collegamento.
    \item Tre: A chiede il MAC di \texttt{192.168.1.1} (il suo gateway) nella subnet 1; R1 chiede quello di \texttt{192.168.2.2} (il
          prossimo hop) nella subnet 2; R2 chiede quello di \texttt{192.168.3.20} (F) nella subnet 3. Ogni richiesta insegna anche al
          destinatario l'associazione di chi chiede.
\end{enumerate}
\end{esercizio}

\begin{esercizio}{Self-learning}
Sei nodi, da A a F, sono collegati a stella a uno switch con la tabella vuota. In ordine: B invia a E, E risponde a B, A invia a B, B
risponde ad A. Dopo ogni evento, mostrare la tabella dello switch e i collegamenti su cui il frame viene inviato, giustificando.

\svolgimento
\begin{enumerate}
    \item B $\rightarrow$ E: lo switch impara B; E è sconosciuto, inonda verso A, C, D, E, F. Tabella: \{B\}.
    \item E $\rightarrow$ B: impara E; B è noto, inoltra solo verso B. Tabella: \{B, E\}.
    \item A $\rightarrow$ B: impara A; inoltra solo verso B. Tabella: \{A, B, E\}.
    \item B $\rightarrow$ A: B è già noto (si aggiorna l'ora); inoltra solo verso A. Tabella: \{A, B, E\}.
\end{enumerate}
\end{esercizio}

\begin{esercizio}{Switch e hub}
Uno switch centrale unisce tre switch di dipartimento, un server Web, un server di posta e un router verso Internet; ogni switch di
dipartimento ha tre host. Ogni collegamento è da 1 Gb/s.
\begin{enumerate}[label=(\alph*)]
    \item Qual è il massimo throughput totale tra i nove host e i due server?
    \item E se i tre switch di dipartimento diventano hub?
    \item E se ogni switch diventa un hub?
\end{enumerate}

\svolgimento
(a) 11 Gb/s: ognuno degli 11 sistemi trasmette a 1 Gb/s nello stesso momento (ogni host verso un altro host del proprio dipartimento, i due
server tra loro), su collegamenti full-duplex senza collisioni. (b) Ogni dipartimento diventa un dominio di collisione da 1 Gb/s
condiviso: $3 + 2 = 5$ Gb/s. (c) Un unico dominio di collisione: 1 Gb/s.
\end{esercizio}

\begin{esercizio}{VLAN}
\begin{enumerate}[label=(\alph*)]
    \item Quante VLAN può configurare al massimo uno switch che supporta 802.1Q? Che cosa fissa il limite?
    \item $N$ switch, ciascuno con gli stessi $K$ gruppi VLAN, devono essere collegati. Quante porte servono senza trunking, e quante con?
\end{enumerate}

\svolgimento
(a) 4094: il VLAN ID ha 12 bit, e i valori 0 e 4095 sono riservati. (b) Per collegare $N$ switch servono $N - 1$ collegamenti (un albero).
Senza trunking ogni collegamento va ripetuto per ciascuna VLAN: $K(N-1)$ cavi, cioè $2K(N-1)$ porte. Con il trunking un cavo porta tutte le
VLAN: $2(N-1)$ porte.
\end{esercizio}

\begin{esercizio}{ARP, visto da ping}
Si cattura con filtro \texttt{arp || icmp} mentre si esegue, uno alla volta:
\begin{lstlisting}[language=bash]
$ ip neigh ; sudo ip neigh flush all
$ ping -c 1 <il proprio gateway>
$ ping -c 1 <un indirizzo della propria subnet che nessuno usa>      # poi: ip neigh
$ ping -c 1 8.8.8.8
\end{lstlisting}
\begin{enumerate}[label=(\alph*)]
    \item Per il gateway: quali frame, in che ordine? Quali campi della richiesta ARP sono ancora sconosciuti, e a chi è indirizzata?
    \item Per l'indirizzo inutilizzato: che cosa va sul filo, quante volte? Che cosa stampa ping, quale macchina ha generato quel messaggio,
          e che stato mostra \texttt{ip neigh}?
    \item Per \texttt{8.8.8.8}: confrontare il MAC di destinazione della richiesta con la riga del gateway in \texttt{ip neigh}. È mai stato
          chiesto \texttt{8.8.8.8}?
\end{enumerate}

\svolgimento
\begin{enumerate}[label=(\alph*)]
    \item ARP request, ARP reply, echo request, echo reply. Nella richiesta il MAC di destinazione (il \emph{target}) è sconosciuto, tutti
          zeri, e il frame è indirizzato in broadcast a \texttt{ff:ff:ff:ff:ff:ff}; la risposta torna in unicast.
    \item Solo richieste ARP, ripetute (tipicamente tre, una al secondo), senza risposta, e nessun echo request. ping stampa \emph{Destination
          Host Unreachable}, generato dalla \textbf{propria} macchina (il suo kernel), e \texttt{ip neigh} mostra l'indirizzo in stato
          \texttt{FAILED} (o \texttt{INCOMPLETE}).
    \item Il MAC di destinazione è quello del gateway. \texttt{8.8.8.8} non viene mai chiesto ad ARP: è fuori dalla subnet, e si risolve solo
          il prossimo hop.
\end{enumerate}
\end{esercizio}

% ══════════════════════════════════════════════════════════════
\section{Livello fisico}
% ══════════════════════════════════════════════════════════════

\begin{esercizio}{I limiti di una linea telefonica}
Un canale telefonico porta le frequenze da 300 a 3400 Hz, con un rapporto segnale/rumore di 30 dB.
\begin{enumerate}[label=(\alph*)]
    \item Qual è la sua banda in hertz? Quanto vale $S/N$ come rapporto puro?
    \item Qual è la sua capacità di Shannon?
    \item Che cosa permette Nyquist con 2 livelli? Con 16?
    \item Un modem invia 2400 simboli al secondo con QAM-16. Qual è la sua velocità in bit? E con 64 punti nella costellazione?
    \item Confrontare (b) con i 33,6 kb/s a cui i modem si sono fermati. Che cosa servirebbe per andare più veloci sulla stessa linea?
\end{enumerate}

\svolgimento
\begin{enumerate}[label=(\alph*)]
    \item $B = 3100$ Hz; $S/N = 10^3 = 1000$.
    \item $C = 3100 \log_2 1001 \approx 3100 \cdot 9{,}97 \approx 30{,}9$ kb/s.
    \item $2B \log_2 V$: 6200 b/s con 2 livelli, 24\,800 b/s con 16.
    \item $2400 \cdot 4 = 9600$ b/s; con 64 punti $2400 \cdot 6 = 14\,400$ b/s.
    \item 33,6 kb/s sta sopra 30,9: le linee reali avevano un $S/N$ un po' migliore di 30 dB (a 35 dB il tetto è circa 36 kb/s), ma il
          tetto c'è comunque. Per andare più veloci servono più banda (togliere il filtro, come fa l'ADSL) o meno rumore (eliminare un tratto
          analogico, come fanno i modem a 56k).
\end{enumerate}
\end{esercizio}

\begin{esercizio}{Codici di linea}
Si prendono i bit \texttt{1 0 1 1 0 0 0 1}.
\begin{enumerate}[label=(\alph*)]
    \item Disegnarli in NRZ; in NRZI, dove un 1 è una transizione e uno 0 nessuna; in Manchester, con un 1 come transizione dal basso
          all'alto a metà bit e uno 0 dall'alto al basso; in AMI, dove uno 0 è zero e gli 1 si alternano, il primo positivo.
    \item Quale dei quattro dà ancora al ricevente una transizione durante una sequenza di otto 0? E di otto 1?
    \item In 4B/5B \texttt{0000} diventa \texttt{11110} e \texttt{0001} diventa \texttt{01001}. Codificare \texttt{0000 0001 0000}. Qual è la
          sequenza di 0 più lunga sulla linea? Quale velocità serve sulla linea per portare 100 Mbit/s di dati?
\end{enumerate}

\svolgimento
\begin{center}
\begin{tikzpicture}[every node/.style={font=\scriptsize\sffamily}]
  \draw[black!20, dashed] (0.00,0.7) -- (0.00,-3.95);
  \draw[black!20, dashed] (0.90,0.7) -- (0.90,-3.95);
  \draw[black!20, dashed] (1.80,0.7) -- (1.80,-3.95);
  \draw[black!20, dashed] (2.70,0.7) -- (2.70,-3.95);
  \draw[black!20, dashed] (3.60,0.7) -- (3.60,-3.95);
  \draw[black!20, dashed] (4.50,0.7) -- (4.50,-3.95);
  \draw[black!20, dashed] (5.40,0.7) -- (5.40,-3.95);
  \draw[black!20, dashed] (6.30,0.7) -- (6.30,-3.95);
  \draw[black!20, dashed] (7.20,0.7) -- (7.20,-3.95);
  \node[font=\scriptsize\ttfamily] at (0.45,0.6) {1};
  \node[font=\scriptsize\ttfamily] at (1.35,0.6) {0};
  \node[font=\scriptsize\ttfamily] at (2.25,0.6) {1};
  \node[font=\scriptsize\ttfamily] at (3.15,0.6) {1};
  \node[font=\scriptsize\ttfamily] at (4.05,0.6) {0};
  \node[font=\scriptsize\ttfamily] at (4.95,0.6) {0};
  \node[font=\scriptsize\ttfamily] at (5.85,0.6) {0};
  \node[font=\scriptsize\ttfamily] at (6.75,0.6) {1};
  \node[anchor=east, font=\scriptsize\bfseries\sffamily] at (-0.15,0.00) {NRZ};
  \draw[black!25] (0,0.00) -- (7.20,0.00);
  \draw[thick, netblue] (0.000,0.300) -- (0.450,0.300) -- (0.450,0.300) -- (0.900,0.300) -- (0.900,-0.300) -- (1.350,-0.300) -- (1.350,-0.300) -- (1.800,-0.300) -- (1.800,0.300) -- (2.250,0.300) -- (2.250,0.300) -- (2.700,0.300) -- (2.700,0.300) -- (3.150,0.300) -- (3.150,0.300) -- (3.600,0.300) -- (3.600,-0.300) -- (4.050,-0.300) -- (4.050,-0.300) -- (4.500,-0.300) -- (4.500,-0.300) -- (4.950,-0.300) -- (4.950,-0.300) -- (5.400,-0.300) -- (5.400,-0.300) -- (5.850,-0.300) -- (5.850,-0.300) -- (6.300,-0.300) -- (6.300,0.300) -- (6.750,0.300) -- (6.750,0.300) -- (7.200,0.300);
  \node[anchor=east, font=\scriptsize\bfseries\sffamily] at (-0.15,-1.10) {NRZI};
  \draw[black!25] (0,-1.10) -- (7.20,-1.10);
  \draw[thick, netblue] (0.000,-0.800) -- (0.450,-0.800) -- (0.450,-0.800) -- (0.900,-0.800) -- (0.900,-0.800) -- (1.350,-0.800) -- (1.350,-0.800) -- (1.800,-0.800) -- (1.800,-1.400) -- (2.250,-1.400) -- (2.250,-1.400) -- (2.700,-1.400) -- (2.700,-0.800) -- (3.150,-0.800) -- (3.150,-0.800) -- (3.600,-0.800) -- (3.600,-0.800) -- (4.050,-0.800) -- (4.050,-0.800) -- (4.500,-0.800) -- (4.500,-0.800) -- (4.950,-0.800) -- (4.950,-0.800) -- (5.400,-0.800) -- (5.400,-0.800) -- (5.850,-0.800) -- (5.850,-0.800) -- (6.300,-0.800) -- (6.300,-1.400) -- (6.750,-1.400) -- (6.750,-1.400) -- (7.200,-1.400);
  \node[anchor=east, font=\scriptsize\bfseries\sffamily] at (-0.15,-2.20) {Manchester};
  \draw[black!25] (0,-2.20) -- (7.20,-2.20);
  \draw[thick, netblue] (0.000,-2.500) -- (0.450,-2.500) -- (0.450,-1.900) -- (0.900,-1.900) -- (0.900,-1.900) -- (1.350,-1.900) -- (1.350,-2.500) -- (1.800,-2.500) -- (1.800,-2.500) -- (2.250,-2.500) -- (2.250,-1.900) -- (2.700,-1.900) -- (2.700,-2.500) -- (3.150,-2.500) -- (3.150,-1.900) -- (3.600,-1.900) -- (3.600,-1.900) -- (4.050,-1.900) -- (4.050,-2.500) -- (4.500,-2.500) -- (4.500,-1.900) -- (4.950,-1.900) -- (4.950,-2.500) -- (5.400,-2.500) -- (5.400,-1.900) -- (5.850,-1.900) -- (5.850,-2.500) -- (6.300,-2.500) -- (6.300,-2.500) -- (6.750,-2.500) -- (6.750,-1.900) -- (7.200,-1.900);
  \node[anchor=east, font=\scriptsize\bfseries\sffamily] at (-0.15,-3.30) {AMI};
  \draw[black!25] (0,-3.30) -- (7.20,-3.30);
  \draw[thick, netblue] (0.000,-3.000) -- (0.450,-3.000) -- (0.450,-3.000) -- (0.900,-3.000) -- (0.900,-3.300) -- (1.350,-3.300) -- (1.350,-3.300) -- (1.800,-3.300) -- (1.800,-3.600) -- (2.250,-3.600) -- (2.250,-3.600) -- (2.700,-3.600) -- (2.700,-3.000) -- (3.150,-3.000) -- (3.150,-3.000) -- (3.600,-3.000) -- (3.600,-3.300) -- (4.050,-3.300) -- (4.050,-3.300) -- (4.500,-3.300) -- (4.500,-3.300) -- (4.950,-3.300) -- (4.950,-3.300) -- (5.400,-3.300) -- (5.400,-3.300) -- (5.850,-3.300) -- (5.850,-3.300) -- (6.300,-3.300) -- (6.300,-3.600) -- (6.750,-3.600) -- (6.750,-3.600) -- (7.200,-3.600);
\end{tikzpicture}
\end{center}
\begin{enumerate}[label=(\alph*)]
    \item[] (a) Nella figura, il livello iniziale di NRZI è basso.
    \item[] (b) Otto 0: solo \textbf{Manchester} (una transizione in ogni bit); NRZ, NRZI e AMI restano piatti. Otto 1: Manchester, NRZI (una
          transizione per ogni 1) e AMI (gli 1 si alternano); NRZ resta piatto.
    \item[] (c) \texttt{11110\,01001\,11110}: la sequenza di 0 più lunga è di \textbf{2}. Servono 125 Mbaud sulla linea (5 bit di segnale ogni
          4 di dati).
\end{enumerate}
\end{esercizio}

% ══════════════════════════════════════════════════════════════
\section{Sicurezza}
% ══════════════════════════════════════════════════════════════

\begin{esercizio}{Quante chiavi}
$N$ persone vogliono parlare ciascuna in privato con ciascuna delle altre $N - 1$; tutto ciò che due si dicono è visibile a tutti. Quante
chiavi servono con la crittografia simmetrica? E con le chiavi pubbliche?

\svolgimento
Simmetrica: una chiave per coppia, $N(N-1)/2$. Pubblica: una coppia per persona, $N$ coppie ($2N$ chiavi), e le chiavi pubbliche si possono
pubblicare.
\end{esercizio}

\begin{esercizio}{Sostituzione e trasposizione}
Un cifrario per sostituzione associa l'alfabeto \texttt{abcdefghijklmnopqrstuvwxyz} a \texttt{QWERTYUIOPASDFGHJKLZXCVBNM}.
\begin{enumerate}[label=(\alph*)]
    \item Cifrare \texttt{this is an easy problem}; decifrare \texttt{ATNL DXLZ LZQN LTEKTZ}.
    \item Con la chiave MEGABUCK, scrivere \texttt{meetmeatthestationatnoon} in righe di otto sotto la chiave, e leggere le colonne
          nell'ordine alfabetico delle lettere della chiave. Qual è il testo cifrato?
    \item Dato solo un testo cifrato, come si capisce quale dei due metodi l'ha prodotto?
\end{enumerate}

\svolgimento
\begin{enumerate}[label=(\alph*)]
    \item \texttt{ZIOL OL QF TQLN HKGWSTD}; \texttt{keys must stay secret}.
    \item Righe \texttt{meetmeat}, \texttt{thestati}, \texttt{onatnoon}. L'ordine delle colonne è A, B, C, E, G, K, M, U, cioè le colonne 4,
          5, 7, 2, 3, 8, 1, 6: \texttt{TSTMTNATOEHNEEATINMTOEAO}.
    \item Dalle frequenze delle lettere: una trasposizione conserva quelle della lingua (le lettere più frequenti restano \emph{e}, \emph{t},
          \emph{a}, \dots), una sostituzione le sposta su altre lettere.
\end{enumerate}
\end{esercizio}

\begin{esercizio}{Un cifrario a blocchi di 3 bit, con e senza CBC}
Un cifrario a blocchi lavora su 3 bit con questa tabella:
\begin{center}
\small\ttfamily
\begin{tabular}{l*{8}{c}}
\toprule
\textrm{Ingresso} & 000 & 001 & 010 & 011 & 100 & 101 & 110 & 111 \\
\textrm{Uscita} & 110 & 111 & 101 & 100 & 011 & 010 & 000 & 001 \\
\bottomrule
\end{tabular}
\end{center}
Il testo in chiaro è \texttt{100100100}.
\begin{enumerate}[label=(\alph*)]
    \item Senza CBC, qual è il testo cifrato?
    \item Trudy lo intercetta. Sa che si usa un cifrario a blocchi di 3 bit senza CBC, ma non la tabella. Che cosa può concludere?
    \item Con CBC e IV $= \texttt{111}$, qual è il testo cifrato? Che cosa deve sapere il ricevente per decifrarlo?
\end{enumerate}

\svolgimento
(a) \texttt{011 011 011}. (b) Che i tre blocchi di testo in chiaro sono \textbf{uguali}, anche senza sapere quali: ECB rivela le ripetizioni.
(c) $c(1) = K(\texttt{100} \oplus \texttt{111}) = K(\texttt{011}) = \texttt{100}$, $c(2) = K(\texttt{100} \oplus \texttt{100}) =
K(\texttt{000}) = \texttt{110}$, $c(3) = K(\texttt{100} \oplus \texttt{110}) = K(\texttt{010}) = \texttt{101}$: \texttt{100 110 101}. Il
ricevente deve conoscere la chiave (la tabella) e l'IV, che viaggia in chiaro.
\end{esercizio}

\begin{esercizio}{RSA con numeri piccoli}
RSA con $p = 5$ e $q = 11$.
\begin{enumerate}[label=(\alph*)]
    \item Quanto valgono $n$ e $z$?
    \item Perché $e = 3$ è una scelta accettabile?
    \item Trovare un $d$ con $de \equiv 1 \pmod z$ e $d < 160$.
    \item Cifrare $m = 8$, poi decifrare il risultato, usando $[(a \bmod n)(b \bmod n)] \bmod n = ab \bmod n$ per tenere piccoli i numeri.
\end{enumerate}

\svolgimento
(a) $n = 55$, $z = 40$. (b) $3 < n$ e non ha fattori comuni con 40. (c) $d = 27$ ($81 \bmod 40 = 1$); vanno bene anche 67, 107, 147. (d) $c
= 8^3 \bmod 55 = 512 \bmod 55 = 17$. Decifrando, $17^{27} \bmod 55$: $17^2 \bmod 55 = 14$, $17^4 = 14^2 \bmod 55 = 31$, $17^8 = 31^2 \bmod
55 = 26$, $17^{16} = 26^2 \bmod 55 = 16$; poiché $27 = 16 + 8 + 2 + 1$, $16 \cdot 26 \cdot 14 \cdot 17 \bmod 55 = 8$.
\end{esercizio}

\begin{esercizio}{Diffie-Hellman con numeri piccoli}
Diffie-Hellman con $p = 11$ e $g = 2$; Alice sceglie $S_A = 4$, Bob $S_B = 7$.
\begin{enumerate}[label=(\alph*)]
    \item Dimostrare che le due chiavi calcolate da Alice e Bob sono sempre uguali.
    \item Calcolare i valori pubblici $T_A$ e $T_B$.
    \item Calcolare la chiave condivisa.
    \item Disegnare un diagramma temporale, una linea ciascuno per Alice, Bob e Trudy, in cui Trudy sta nel mezzo e finisce per condividere
          una chiave con ciascuno.
\end{enumerate}

\svolgimento
(a) $T_B^{S_A} = (g^{S_B})^{S_A} = g^{S_A S_B} = (g^{S_A})^{S_B} = T_A^{S_B} \pmod p$. (b) $T_A = 2^4 \bmod 11 = 5$, $T_B = 2^7 \bmod 11 =
7$. (c) $7^4 \bmod 11 = 3$ e $5^7 \bmod 11 = 3$: la chiave è 3. (d) Trudy sceglie per esempio $S_T = 3$, quindi $T_T = 8$:
\begin{center}
\begin{tikzpicture}[yscale=0.75, every node/.style={font=\scriptsize\sffamily}]
  \timeline{al}{0}{Alice}{3.6}
  \timeline{tr}{4.5}{Trudy}{3.6}
  \timeline{bo}{9}{Bob}{3.6}
  \draw[msg, netblue] (0,-0.5) -- node[above] {$T_A = 5$} (4.5,-0.9);
  \draw[msg, netred] (4.5,-1.0) -- node[above] {$T_T = 8$} (9,-1.4);
  \draw[msg, netblue] (9,-1.7) -- node[above] {$T_B = 7$} (4.5,-2.1);
  \draw[msg, netred] (4.5,-2.2) -- node[above] {$T_T = 8$} (0,-2.6);
  \node[anchor=east, text width=2.6cm, align=right] at (-0.2,-3.1) {chiave $8^4 \bmod 11 = 4$};
  \node[anchor=west, text width=2.6cm] at (9.2,-3.1) {chiave $8^7 \bmod 11 = 2$};
  \node[text width=3.2cm, align=left, anchor=west, text=netred] at (4.65,-3.3) {con Alice: $5^3 \bmod 11 = 4$\\con Bob: $7^3 \bmod 11 = 2$};
\end{tikzpicture}
\end{center}
\end{esercizio}

\begin{esercizio}{Autenticarsi nei due versi}
Alice si autentica presso Bob così: invia \emph{sono Alice}; Bob risponde con un nonce fresco $R$; Alice restituisce $K_{A\text{-}B}(R)$,
cifrato con la chiave che condivide con Bob; Bob lo decifra e controlla che sia $R$. Come visto nel corso, funziona.
\begin{enumerate}[label=(\alph*)]
    \item Ora anche Bob deve autenticarsi presso Alice, con lo stesso protocollo e la stessa chiave: invia \emph{sono Bob}, lei risponde con
          un nonce, lui lo restituisce cifrato.
    \item Le due esecuzioni possono intrecciarsi in qualsiasi ordine. Mostrare come Trudy, che non conosce $K_{A\text{-}B}$, può far credere a
          Bob di essere Alice. Suggerimento: far rispondere Bob alla sua stessa domanda.
    \item Proporre una modifica che fermi l'attacco.
\end{enumerate}

\svolgimento
\begin{enumerate}[label=(\alph*)]
    \item Bob: \emph{sono Bob}; Alice: $R'$; Bob: $K_{A\text{-}B}(R')$; Alice decifra e controlla.
    \item Trudy invia a Bob \emph{sono Alice}, e Bob risponde con il nonce $R$. Trudy non sa calcolare $K_{A\text{-}B}(R)$. Ma quando Bob
          comincia la propria autenticazione presso ``Alice'' (\emph{sono Bob}), Trudy risponde con lo stesso $R$ come se fosse il suo nonce;
          Bob restituisce $K_{A\text{-}B}(R)$, e Trudy lo usa nella prima esecuzione: Bob ha risposto alla sua stessa domanda, e accetta Trudy
          come Alice (l'attacco per riflessione).
    \item Chiavi diverse per i due versi ($K_{A \rightarrow B} \neq K_{B \rightarrow A}$); oppure includere l'identità nella risposta cifrata,
          $K_{A\text{-}B}(R, \text{Alice})$; oppure Bob rifiuta di rispondere a un nonce che lui stesso ha emesso ed è ancora in sospeso.
\end{enumerate}
\end{esercizio}

\begin{esercizio}{TLS: quattro domande veloci}
Rispondere con l'handshake e i record di TLS visti nel corso.
\begin{enumerate}[label=(\alph*)]
    \item Vero o falso: un record TLS ha un campo per i numeri di sequenza TLS. Dove va allora il numero di sequenza?
    \item A che cosa servono i nonce casuali nell'handshake TLS?
    \item Una sessione TLS usa un cifrario a blocchi con CBC. Vero o falso: il server invia l'IV al client in chiaro.
    \item Bob apre una connessione TCP verso Trudy, che finge di essere Alice e invia a Bob il certificato di Alice. In quale passo
          dell'handshake Bob scopre di non parlare con Alice?
\end{enumerate}

\svolgimento
\begin{enumerate}[label=(\alph*)]
    \item Falso: il numero di sequenza non viaggia. Entrambe le parti lo contano, e lo includono nel calcolo dell'HMAC di ogni record.
    \item A rendere unica ogni sessione: le chiavi dipendono dai due nonce, quindi una sessione registrata e ripetuta il giorno dopo produce
          chiavi diverse, e i record ripetuti falliscono l'HMAC.
    \item Vero in TLS 1.1 e 1.2: ogni record CBC porta in testa un IV esplicito, in chiaro (in TLS 1.0, e nel ``quasi-TLS'' del Kurose, gli
          IV si derivano invece dal master secret).
    \item Trudy non ha la chiave privata di Alice, quindi non sa aprire il pre-master secret che Bob ha cifrato con la chiave pubblica di
          Alice, e non sa calcolare le chiavi: il messaggio con l'HMAC dell'handshake (\emph{Finished}) che Trudy dovrebbe inviare non torna,
          e Bob se ne accorge lì (passo 6).
\end{enumerate}
\end{esercizio}

\begin{esercizio}{Un handshake TLS da leggere}
Si cattura con filtro \texttt{tls}:
\begin{lstlisting}[language=bash]
$ curl -s -o /dev/null --tls-max 1.2 https://www.unina.it
$ curl -s -o /dev/null https://www.unina.it
\end{lstlisting}
Il primo forza TLS 1.2, il cui handshake è quello seguito nel corso; il secondo permette l'1.3, per vedere che cosa nasconde la versione
più nuova. Il server può rispondere 403 a curl: non importa, l'handshake avviene prima che il server veda la richiesta.
\begin{enumerate}[label=(\alph*)]
    \item Nella cattura 1.2, elencare i messaggi dell'handshake in ordine. Dove sono i due nonce?
    \item Quale sito si sta visitando è leggibile in chiaro: in quale messaggio, quale campo?
    \item Aprire il certificato: soggetto, emittente, validità. Chi garantisce per chi?
    \item Da quale record in poi tutto è cifrato?
    \item Nella cattura 1.3, quali di queste cose si leggono ancora? Che cosa ha perso un osservatore di rete?
\end{enumerate}

\svolgimento
\begin{enumerate}[label=(\alph*)]
    \item Client Hello; Server Hello, Certificate, Server Key Exchange, Server Hello Done; Client Key Exchange, Change Cipher Spec, Finished;
          Change Cipher Spec, Finished; poi Application Data. I nonce sono i campi \texttt{Random} del Client Hello e del Server Hello.
    \item Nel Client Hello, estensione \texttt{server\_name} (la SNI).
    \item Il soggetto è \texttt{www.unina.it} (o un nome che lo comprende); l'emittente è una CA intermedia, il cui certificato è firmato a
          sua volta da una radice presente nel sistema: la radice garantisce per l'intermedia, l'intermedia per il sito.
    \item Dal primo \emph{Finished}, che segue il Change Cipher Spec: da lì in poi tutto è cifrato.
    \item In 1.3 si leggono solo Client Hello (con la SNI) e Server Hello; certificato, Finished e tutto il resto viaggiano come Application
          Data. Un osservatore perde la vista del certificato (e quindi della catena e del soggetto): resta la SNI, finché non c'è ECH.
\end{enumerate}
\end{esercizio}

\begin{esercizio}{Un firewall stateful}
La rete interna è \texttt{222.22.0.0/16}. Scrivere la tabella dei filtri e la tabella delle connessioni di un firewall stateful, il più
restrittivo possibile, che (a) permetta a ogni utente interno di aprire sessioni Telnet verso host esterni; (b) permetta agli utenti esterni
di navigare il sito aziendale su \texttt{222.22.0.12}; (c) blocchi ogni altro pacchetto in entrata e in uscita. La tabella delle connessioni
contiene ora tre connessioni, tutte aperte dall'interno. Inventare indirizzi e porte, e dire per ogni riga della tabella dei filtri se
consulta la tabella delle connessioni.

\svolgimento
\begin{center}
\small
\begin{tabular}{lllcccc c}
\toprule
\textbf{Azione} & \textbf{Sorgente} & \textbf{Destinazione} & \textbf{Prot.} & \textbf{P. sorg.} & \textbf{P. dest.} & \textbf{Flag} &
\textbf{Conn.} \\
\midrule
allow & \texttt{222.22/16} & esterno & TCP & $>1023$ & 23 & qualsiasi & \\
allow & esterno & \texttt{222.22/16} & TCP & 23 & $>1023$ & ACK & X \\
allow & esterno & \texttt{222.22.0.12} & TCP & $>1023$ & 80 & qualsiasi & \\
allow & \texttt{222.22.0.12} & esterno & TCP & 80 & $>1023$ & ACK & X \\
deny & tutto & tutto & tutto & tutto & tutto & tutto & \\
\bottomrule
\end{tabular}
\end{center}
Le righe marcate X lasciano passare un pacchetto solo se appartiene a una connessione nella tabella: le risposte Telnet solo per sessioni
aperte dall'interno, le risposte del server Web solo per connessioni aperte con un SYN passato dalla terza riga. Tabella delle connessioni:
\begin{center}
\small\ttfamily
\begin{tabular}{llll}
\toprule
\textrm{\textbf{Sorgente}} & \textrm{\textbf{Destinazione}} & \textrm{\textbf{P. sorg.}} & \textrm{\textbf{P. dest.}} \\
\midrule
222.22.1.7 & 37.96.87.123 & 12699 & 23 \\
222.22.93.2 & 199.1.205.23 & 37654 & 23 \\
222.22.65.143 & 203.77.240.43 & 48712 & 23 \\
\bottomrule
\end{tabular}
\end{center}
\end{esercizio}

% ══════════════════════════════════════════════════════════════
\section{Anonimato}
% ══════════════════════════════════════════════════════════════

\begin{esercizio}{Chi sa che cosa in un circuito}
Alice raggiunge il sito S attraverso Tor, su un circuito con un guard G, un relay intermedio M e un exit X.
\begin{enumerate}[label=(\alph*)]
    \item Per G, M, X e S, dire se conosce l'indirizzo IP di Alice, il nome di S e il contenuto della richiesta. Rispondere due volte: con
          S raggiunto in \texttt{http://} e in \texttt{https://}.
    \item Un'osservatrice vede sia il traffico tra Alice e G sia quello tra X e S, ma non può decifrare nulla. Che cosa può ancora imparare,
          e come?
    \item Perché tre relay? Che cosa rivelerebbe un relay solo, e che cosa due?
\end{enumerate}

\svolgimento
\begin{enumerate}[label=(\alph*)]
    \item \begin{tabular}[t]{lccc}
          \toprule
           & \textbf{IP di Alice} & \textbf{nome di S} & \textbf{contenuto} \\
          \midrule
          G & sì & no & no \\
          M & no & no & no \\
          X & no & sì & sì in \texttt{http}, no in \texttt{https} \\
          S & no (vede X) & sì & sì \\
          \bottomrule
          \end{tabular}

          In \texttt{https} l'exit vede solo nome (SNI) e indirizzo di S; in entrambi i casi G non vede né S né il contenuto, cifrati dagli
          strati di M e X.
    \item Che Alice sta parlando con S: confronta tempi, raffiche e volumi del traffico ai due capi (correlazione temporale). La cifratura non
          cambia quando i byte si muovono né quanti sono.
    \item Un relay conosce entrambi gli estremi. Con due, l'exit conosce l'ingresso, che diventa il nodo da sorvegliare. Con tre l'exit
          conosce solo un intermedio, e l'intermedio nessuno dei due estremi.
\end{enumerate}
\end{esercizio}

% ══════════════════════════════════════════════════════════════
\section{Tutta la pila}
% ══════════════════════════════════════════════════════════════

\begin{esercizio}{Tutta la pila, su carta}
Si entra in una stanza, si collega il portatile all'Ethernet e si scarica una pagina Web. Nulla è in cache: nessun indirizzo, nessuna riga
ARP, nessuna risposta DNS.
\begin{enumerate}[label=(\alph*)]
    \item Elencare ogni passo di protocollo dall'accensione del portatile alla visualizzazione della pagina: Ethernet, DHCP, ARP, DNS, TCP,
          HTTP.
    \item Dire esplicitamente come il portatile impara l'indirizzo IP del proprio gateway, e poi il suo indirizzo MAC.
    \item Per ogni passo, quali indirizzi sono noti e quali ancora sconosciuti?
\end{enumerate}

\svolgimento
\begin{enumerate}[label=(\alph*)]
    \item DHCP Discover (da \texttt{0.0.0.0:68} a \texttt{255.255.255.255:67}, frame in broadcast), Offer, Request, ACK; ARP per il gateway
          (o per il server DNS, se è sulla LAN); query DNS per il nome del sito (UDP 53, verso il server DNS attraverso il gateway) e
          risposta; handshake TCP con il server Web (SYN, SYN-ACK, ACK); GET HTTP e risposta; chiusura della connessione.
    \item Il DHCP ACK porta, oltre all'indirizzo e alla maschera, l'indirizzo IP del gateway (opzione \emph{router}) e quello del server
          DNS. Il MAC del gateway lo dà poi una richiesta ARP in broadcast per quell'indirizzo.
    \item Al DHCP il portatile conosce solo il proprio MAC (usa \texttt{0.0.0.0} e i broadcast). Dopo l'ACK conosce il proprio IP, la maschera,
          l'IP del gateway e del DNS, ma non il MAC del gateway; dopo l'ARP anche quello. Per la query DNS conosce l'IP del server DNS ma non
          quello del sito; dopo la risposta, l'IP del sito. Il MAC del sito non lo conoscerà mai: ogni frame va al gateway.
\end{enumerate}
\end{esercizio}

\begin{esercizio}{La vostra chat, tutta la pila}
Server su una macchina, client su un'altra della stessa LAN; cattura sull'interfaccia reale del client, senza filtro. Sul client, in ordine:
\begin{lstlisting}[language=bash]
$ nmcli connection down <conn> ; nmcli connection up <conn>    # solo sulla propria macchina
$ sudo ip neigh flush all ; resolvectl flush-caches
$ ./chat <nome o indirizzo del server> 5000                     # un messaggio breve, poi quit
\end{lstlisting}
Si rinnova l'indirizzo solo sulla propria macchina: in laboratorio si salta quella riga. Ci si collega per nome se il server ne ha uno nel
DNS; altrimenti per indirizzo, e la parte DNS si salta.
\begin{enumerate}[label=(\alph*)]
    \item Costruire la sequenza temporale: ogni protocollo che compare, in ordine, con il numero del suo primo pacchetto. Quanti messaggi
          DHCP si vedono, e perché forse non quattro?
    \item Per ogni scambio ARP e DNS: chi ha chiesto, su chi, e perché in quel momento?
    \item Si prende il frame che porta il messaggio breve. Livello per livello, dare la lunghezza di ogni header e il campo che dice che cosa
          viene dopo: EtherType, protocollo IP, porta di destinazione, e il proprio byte di tipo. Verificare che le lunghezze tornino con la
          lunghezza del frame.
    \item Inviare un messaggio con un corpo di 3000 byte, prima con la versione connessa della chat, poi con quella a datagrammi (alzando
          la lunghezza massima dichiarata, se serve). Quanti frame produce ciascuna? Quale livello ha tagliato il messaggio, e come lo
          rimette insieme il ricevente?
    \item Nella versione a datagrammi, che cosa succede all'intero messaggio se uno di quei frame si perde? E in quella connessa?
    \item Disegnare il diagramma temporale completo di (a), una linea verticale per macchina coinvolta: client, server, gateway, server DHCP e
          DNS.
\end{enumerate}

\svolgimento
\begin{enumerate}[label=(\alph*)]
    \item DHCP, ARP (spesso un ARP ``gratuito'' o di verifica del proprio indirizzo, poi la richiesta per il gateway), DNS, ARP per il
          server della chat (è sulla stessa LAN), TCP (handshake), i dati della chat, la chiusura. Al rinnovo di un indirizzo già noto il
          client salta la scoperta: si vedono solo \textbf{Request e ACK}, due messaggi invece di quattro.
    \item ARP per il gateway (o per il server DNS) prima della query DNS, perché il frame della query ha bisogno di un MAC di destinazione;
          DNS dal client verso il resolver, per il nome del server, prima della \texttt{connect}; ARP per il server della chat, chiesto dal
          client subito prima del SYN.
    \item Ethernet 14 byte (EtherType \texttt{0x0800}), IP 20 (protocollo 6), TCP 32 con le opzioni timestamp (porta di destinazione 5000),
          poi il nostro header di 5 byte (quattro di lunghezza e il byte di tipo) e il corpo. Per un corpo di 5 byte: $14 + 20 + 32 + 5 + 5 =
          76$ byte, la lunghezza stampata da Wireshark (senza preambolo e CRC).
    \item Connessa: TCP taglia il flusso in segmenti di al più 1448 byte, quindi \textbf{3 frame}; il ricevente riceve un flusso, e il nostro
          framing (la lunghezza in testa) ricompone il messaggio. A datagrammi: un datagramma UDP di 3013 byte, che IP \textbf{frammenta} in 3
          frammenti; il riassemblaggio lo fa IP del ricevente prima di consegnare il datagramma.
    \item A datagrammi, se si perde un frammento si perde l'intero datagramma, e quindi l'intero messaggio (a meno che la chat non lo
          ritrasmetta). Nella versione connessa TCP ritrasmette solo il segmento perso, e il messaggio arriva intero.
    \item Client, server DHCP (spesso il gateway stesso), gateway, server DNS e server della chat: DHCP tra client e server DHCP; ARP tra
          client e gateway; DNS tra client e resolver attraverso il gateway; ARP tra client e server della chat; handshake, dati e chiusura
          tra client e server.
\end{enumerate}
\end{esercizio}

\begin{esercizio}{Quanto costa la vostra chat sotto carico}
Si prendono le versioni concorrenti della chat (con processi, con thread, a eventi e con un pool) e le si misura sotto lo stesso carico.
\begin{itemize}
    \item Il carico: un client che apre $N$ connessioni; ciascuna invia un messaggio ogni 100 ms, con l'istante d'invio, e registra quanto
          impiega ogni messaggio a raggiungere un altro client. Si esegue per 60 secondi.
    \item Le misure, per $N = 10$, 50 e 200: media e 95º percentile del ritardo; messaggi consegnati al secondo; tempo di CPU e memoria
          residente del server (\texttt{ps -o rss,time}).
    \item Tutto in una tabella: una riga per versione, un gruppo di colonne per $N$.
\end{itemize}
\begin{enumerate}[label=(\alph*)]
    \item Quale versione peggiora per prima al crescere di $N$, e qual è il suo collo di bottiglia: CPU, un lock, una coda per client, il
          numero di thread? Mostrarlo dai numeri, non dalle lezioni.
    \item Durante un'esecuzione con $N = 50$, sospendere un client con Ctrl-Z per dieci secondi. Quali versioni rallentano tutti, e perché
          proprio quelle?
    \item Server e client su due macchine. A quale $N$ la rete diventa il collo di bottiglia, e come lo si capisce dalle misure? Confrontare
          il tasso consegnato con la velocità del collegamento.
    \item Con 200 client, quanti byte dei messaggi sono header di TCP, IP ed Ethernet? La versione a datagrammi lo cambierebbe?
\end{enumerate}

\svolgimento
È un esercizio di misura: i numeri sono i vostri. Che cosa aspettarsi e come leggerli:
\begin{enumerate}[label=(\alph*)]
    \item La chat è un \emph{broadcast}: ogni messaggio va a tutti, quindi il lavoro cresce come $N^2$ (a $N = 200$, 2000 messaggi al secondo
          da consegnare a 199 destinatari). La versione con processi paga i passaggi tra processi e la memoria (un processo per client: la RSS
          cresce linearmente), quella con thread un \textbf{lock} sulla lista dei client tenuto durante l'invio a tutti, quella a eventi
          l'unico thread che fa tutto (CPU al 100\% su un core). Il collo di bottiglia si legge dal 95º percentile che esplode e dalla CPU o
          dalla memoria che lo accompagnano.
    \item Rallentano tutti le versioni che inviano a ciascun client con \texttt{write} \textbf{bloccanti} tenendo un lock, o dall'unico
          thread: quando il buffer del client sospeso si riempie, la \texttt{write} verso di lui si blocca e ferma la consegna agli altri. Le
          versioni con una coda per client (o con scritture non bloccanti) non ne risentono.
    \item Quando il tasso consegnato smette di crescere con $N$ mentre la CPU del server resta sotto il massimo, e si avvicina alla velocità
          del collegamento (contando gli header): il collo di bottiglia è la rete.
    \item Per un messaggio breve, $14 + 20 + 32 = 66$ byte di header (più preambolo, CRC e intervallo tra frame sul filo) contro una decina di
          byte di dati: la maggior parte dei byte sono header. Con UDP l'header di trasporto scende da 32 a 8 byte, ma resta la stessa
          proporzione di fondo, e si perde l'affidabilità.
\end{enumerate}
\end{esercizio}
